% defs.tex — data-driven macros, filled from final measurements (results/*.json).
% --- Table I: cold-wake energy (J) and cold-wake-in-inferences per model ---
\newcommand{\CEsq}{0.34}\newcommand{\CIsq}{4}
\newcommand{\CEsh}{0.85}\newcommand{\CIsh}{19}
\newcommand{\CEmb}{1.21}\newcommand{\CImb}{11}
\newcommand{\CEgo}{2.39}\newcommand{\CIgo}{7}
\newcommand{\CEde}{3.82}\newcommand{\CIde}{5}
\newcommand{\CErEighteen}{4.01}\newcommand{\CIrEighteen}{12}
\newcommand{\CEef}{3.95}\newcommand{\CIef}{10}
\newcommand{\CErFifty}{8.80}\newcommand{\CIrFifty}{12}
\newcommand{\CEvit}{35.2}\newcommand{\CIvit}{12}
\newcommand{\CEsqi}{0.31}\newcommand{\CIsqi}{4}
\newcommand{\CEmbi}{0.47}\newcommand{\CImbi}{11}
\newcommand{\CErFiftyi}{2.98}\newcommand{\CIrFiftyi}{13}
% --- ViT-Base row (E1) ---
\newcommand{\VITS}{377}\newcommand{\VITW}{4620}\newcommand{\VITT}{12.3}\newcommand{\VITD}{78}
% --- predictive model ---
\newcommand{\BWFLASH}{\SI{96}{MB/s}}
\newcommand{\WAKERSQDISK}{1.00}
\newcommand{\WAKESLOPE}{\SI{13.2}{ms/MB}}
\newcommand{\WAKEINT}{\SI{53}{ms}}
\newcommand{\WAKERSQ}{0.998}
\newcommand{\WAKEMAPE}{12\%}
% --- amortization ---
\newcommand{\BTENLO}{38}
\newcommand{\BTENHI}{119}
\newcommand{\OVERHEADBONE}{80--96\%}
\newcommand{\OVERHEADBONEE}{73--95\%}
\newcommand{\CINFRANGE}{4--19}
% --- storage / dvfs ---
\newcommand{\STORAGEFRAC}{55--81\%}
\newcommand{\RAMPMEAN}{2.5}
% --- policy (abstract) ---
\newcommand{\ABSPOLICYSAVE}{15\%}
% --- prose fragments ---
\newcommand{\CLIFFDESC}{sits where a co-tenant drives free memory below the model's working set:
ResNet-50 stays warm (${\sim}\SI{344}{ms}$) until the co-tenant reaches ${\sim}\SI{1.4}{GB}$
(leaving ${<}\SI{350}{MB}$ free), beyond which session creation thrashes swap and the wake
explodes to over \SI{30}{s}---a ${>}80\times$ jump. Every model tested shows the same cliff}
\newcommand{\ENERGYPARA}{In absolute terms a cold wake costs from \SI{0.3}{J} (SqueezeNet-INT8) to
\SI{8.8}{J} (ResNet-50), and \SI{35}{J} for ViT-Base, scaling with model size; idle draw is
\SI{2.1}{W}. Because loading is I/O-bound the wake draws less instantaneous power than sustained
compute, but its duration makes it costly: for a device that wakes to run a single inference,
\OVERHEADBONEE\ of the energy spent is cold-start overhead rather than useful work.}
\newcommand{\POLICYPARA}{In the seed-averaged simulation, across both workloads and all budgets,
\gdtax{} never underperforms LRU or LFU, and it reduces energy per request by up to
\SI{15}{\percent} (Zipf) and
\SI{11}{\percent} (surveillance) at tight budgets, tracking---and at low-to-mid budgets matching
or beating---an offline Belady hit-rate reference. The gain is largest where cost and popularity
are misaligned (an infrequent but expensive-to-reload detector), exactly the case LRU and LFU
mishandle by evicting the costly model between its accesses.}
\newcommand{\ONDEVICEPARA}{\gdtax{} reduces measured energy per request in three of the four
tested configurations---by up to \SI{14}{\percent} versus LRU (Zipf, \SI{200}{MB} budget) and
\SI{9}{\percent} (surveillance, \SI{320}{MB})---while also lowering mean latency and miss count;
in the tightest surveillance case it trails LRU by \SI{4}{\percent}, within run-to-run variation.
The measured ranking (reload $\gg$ LRU\,$\approx$\,LFU\,$\gtrsim$\,\gdtax{}) matches the
seed-averaged simulation, and the absolute reduction from caching is large: always-reload spends
\SIrange{2.8}{1.4}{J} per request, which residency cuts to as little as \SI{0.61}{J}.}
\newcommand{\OPTPARA}{The crossover is model-dependent: for most models optimization pays for
itself after only a few inferences ($B^{*}<5$), but for graphs whose optimization is expensive it
is substantial---$B^{*}\approx18$ for ResNet-18---so a doorbell-style deployment that runs a
handful of inferences per wake should disable it.}
